% =====================================================================
%  5. DESCRIPCIÓN DE LA AUDITORÍA  <- NÚCLEO DEL TRABAJO
%  Responsable: INTEGRANTE 3
%  ESTADO: REDACTADO (1 de agosto de 2026) — ~2.300 palabras + 2 tablas
%  Fuente principal: Sección C de docs/bitacora_decisiones.md
%  Documento ancla: HRIA de On Common Ground (PDF espejado en q4cdn.com)
%
%  Contiene la Tabla del equipo auditor y la Tabla de criterios de
%  auditoría. Falta insertar la Figura 3 (línea de tiempo): el entorno
%  está listo y comentado al final del archivo.
%
%  ENFOQUE MANTENIDO: esta sección DESCRIBE la auditoría. La evaluación
%  crítica va en la sección 7. Donde aparece un dato incómodo —el
%  financiamiento por la empresa auditada, la no ejecución de la revisión
%  por pares, la ausencia de Sipacapa— se consigna como HECHO y se remite
%  a la sección 7, sin juzgarlo aquí. Ese orden es lo que distingue un
%  informe de auditoría de un alegato.
% =====================================================================

\section{Descripción de la auditoría}
\label{sec:auditoria}

Esta sección reconstruye el proceso de auditoría analizado: quién lo promovió,
quién lo ejecutó, con qué mandato, contra qué criterios y mediante qué método. El
documento examinado es la \hria, elaborada por \auditora\ y publicada en mayo de
2010 \parencite{oncommonground2010hria}.

La exposición se atiene deliberadamente a un registro descriptivo. Varios de los
hechos que se consignan a continuación —que la evaluación fue financiada por la
empresa auditada, que la revisión por pares prevista no llegó a ejecutarse, o que
la participación del municipio de Sipacapa resultó insuficiente— admiten una
lectura crítica que se desarrolla en la sección \ref{sec:critico}. Separar ambos
planos no es una formalidad: una auditoría se evalúa comparando lo que se hizo
con lo que debió hacerse, y para ello es imprescindible establecer primero, sin
valoración, qué fue exactamente lo que se hizo.

% ---------------------------------------------------------------------
\subsection{Antecedentes y origen del encargo}

El rasgo más singular de este proceso es su origen. La evaluación no fue iniciativa
de la empresa ni orden de una autoridad estatal, sino resultado de la presión
ejercida por un grupo de \textbf{accionistas socialmente responsables} sobre la
matriz corporativa. Los promotores fueron Ethical Funds, el First Swedish
National Pension Fund, el Fourth Swedish National Pension Fund, el Public Service
Alliance of Canada Staff Pension Fund y SHARE
\parencite{oncommonground2010hria}.

Esa vía de activación merece atención analítica. En el esquema clásico de control
ambiental expuesto en la sección \ref{sec:marco}, la verificación de la conducta
empresarial procede del Estado, mediante fiscalización, o del mercado de la
certificación, mediante auditorías de tercera parte. Aquí operó un tercer
mecanismo: el capital accionario. Fueron inversionistas institucionales quienes,
en ejercicio de sus políticas de inversión responsable, condicionaron su relación
con la empresa a que esta se sometiera a una evaluación independiente de derechos
humanos. Se trata de un mecanismo de gobernanza privada que no sustituye a la
fiscalización estatal —carece de potestad sancionadora— pero que en este caso
consiguió lo que aquella no había logrado: que la operación fuese examinada por un
tercero externo y que el resultado se hiciera público.

El contexto que hizo posible esa presión estaba ya formado. En 2005 la
organización no gubernamental Madre Selva había presentado una queja ante la
Compliance Advisor Ombudsman de la Corporación Financiera Internacional respecto
del préstamo otorgado al proyecto, procedimiento que concluyó en mayo de 2006
\parencite{cao2006marlin}. El conflicto social en San Miguel Ixtahuacán y
Sipacapa estaba plenamente instalado, la consulta comunitaria de Sipacapa se había
celebrado en 2005, y la revisión independiente del estudio de impacto ambiental
elaborada por \textcite{moran2004review} circulaba desde 2004. La evaluación de
derechos humanos no surgió, por tanto, en un escenario neutro, sino como respuesta
a un cuestionamiento previo y sostenido.

El proceso se formalizó mediante un \textbf{Memorando de Entendimiento} suscrito en
marzo de 2008. Para su gobierno se constituyó un \textbf{Comité Directivo}
(\textit{Steering Committee}) independiente, integrado por Bob Walker, Manfredo
Marroquín y Dave Deisley, con la participación de Bill Brassington. Fue ese comité
—y no la dirección de la empresa— el que en octubre de 2008 seleccionó a la firma
auditora \parencite{oncommonground2010hria}. La interposición de una instancia
independiente entre quien financia y quien audita es un elemento de diseño
institucional deliberado, cuyo alcance y cuyas limitaciones se discuten en la
sección \ref{sec:critico}.

% ---------------------------------------------------------------------
\subsection{Organismo auditor, mandato y financiamiento}

La evaluación fue ejecutada por \auditora, consultora independiente con sede en
Vancouver, Canadá, especializada en evaluación social y relacionamiento
comunitario en el sector minero \parencite{oncommonground2010hria}.

El encargo procedió del Comité Directivo de la evaluación, actuando «en nombre de
Goldcorp». El \textbf{financiamiento corrió íntegramente a cargo de Goldcorp}, es
decir, de la matriz de la empresa auditada. Este dato se consigna aquí como hecho
verificado y sin valoración; su análisis corresponde a la sección
\ref{sec:critico}. Conviene únicamente advertir que la circunstancia no es
excepcional en el ámbito de las evaluaciones de impacto en derechos humanos:
prácticamente todas se financian de ese modo, sencillamente porque no existe un
mecanismo alternativo de financiación consolidado. Lo que distingue a un proceso
de otro no es la ausencia de esa dependencia, sino las salvaguardas que se
articulan frente a ella.

En este caso, tales salvaguardas fueron tres, todas ellas fijadas por escrito en el
Memorando de Entendimiento. La primera es el conjunto de principios rectores del
proceso: \textbf{transparencia, independencia e inclusividad}. La segunda es la
interposición del Comité Directivo como órgano de selección y supervisión. La
tercera es el compromiso de publicación íntegra del informe, incluidos los
hallazgos desfavorables a la empresa; compromiso que efectivamente se cumplió, y
que explica que este trabajo pueda apoyarse en un documento público de más de
doscientas páginas.

A estos principios el propio equipo auditor añadió, en diciembre de 2008, dos
principios éticos adicionales de su propia iniciativa: el \textbf{consentimiento
informado} de las personas entrevistadas y la \textbf{confidencialidad} de sus
declaraciones \parencite{oncommonground2010hria}. La incorporación no es un
detalle menor. En una auditoría desarrollada en un territorio en conflicto, donde
quienes testimonian pueden sufrir represalias, la confidencialidad no es una
cortesía procedimental sino una condición material para que exista testimonio. Se
corresponde, además, con el principio de confidencialidad que la ISO 19011:2018
exige a todo auditor \parencite{iso2018iso19011}.

% ---------------------------------------------------------------------
\subsection{Composición y competencia del equipo auditor}

La ISO 19011:2018 exige que la competencia del equipo auditor se determine y
evalúe en función de los conocimientos y habilidades necesarios para alcanzar los
objetivos de la auditoría \parencite{iso2018iso19011}. Procede, en consecuencia,
examinar si el perfil del equipo cubría las dos dimensiones —social y ambiental—
comprendidas en su alcance. La Tabla \ref{tab:equipo-auditor} resume su
composición.

% --- TABLA: Composición del equipo auditor ---------------------------
\begin{table}[H]
  \centering
  \caption{Composición del equipo de la evaluación de derechos humanos de la Mina Marlin}
  \label{tab:equipo-auditor}
  \begin{small}
  \begin{tabularx}{\textwidth}{@{}L{3.6cm} L{3.4cm} Y@{}}
    \toprule
    \textbf{Integrante} & \textbf{Rol} & \textbf{Perfil y aporte} \\
    \midrule
    Susan Joyce & Líder del equipo &
      Socióloga con más de 25 años de experiencia en aspectos sociales de la
      minería y en derechos humanos. \\
    \addlinespace
    Myriam Cabrera & Equipo principal & Evaluación social y trabajo de campo. \\
    \addlinespace
    Giselle Huamani & Equipo principal &
      Experiencia regional en transformación de conflictos socioambientales. \\
    \addlinespace
    Lloyd Lipsett & Equipo principal &
      Marco jurídico de empresas y derechos humanos. \\
    \addlinespace
    Monica Leonardo Segura & Apoyo & Contexto jurídico guatemalteco. \\
    \addlinespace
    Sandro Macassi & Apoyo & Análisis de conflictos y comunicación. \\
    \addlinespace
    International Alert & Revisión por pares &
      \textbf{Prevista en el diseño, no ejecutada según lo planeado.} \\
    \bottomrule
  \end{tabularx}
  \end{small}
  \par\vspace{3pt}
  \raggedright\footnotesize
  Nota. Los perfiles de los integrantes distintos de la líder del equipo se
  describen de forma resumida; el informe original no detalla la trayectoria
  individual de cada uno. Elaboración propia a partir de
  \textcite{oncommonground2010hria}.
\end{table}

Del cuadro se desprenden dos observaciones de carácter descriptivo. La primera es
que el equipo reunía un perfil predominantemente \textbf{social y jurídico}:
sociología, derechos humanos, análisis de conflictos y derecho. No incorporaba
especialistas en hidrogeología, geoquímica ni ingeniería de relaves. Esta
composición es coherente con el objeto declarado de la evaluación —una evaluación
de impactos en derechos humanos, no una auditoría técnico-ambiental— y el equipo
resolvió la brecha por la vía de encargar una revisión ambiental independiente,
según se detalla en la subsección de metodología.

La segunda observación se refiere a la revisión por pares. El diseño del proceso
contemplaba que International Alert, organización con sede en Londres
especializada en construcción de paz, ejerciera de revisor par del trabajo. Esa
revisión \textbf{no llegó a ejecutarse según lo previsto}. Se consigna aquí como
hecho, por su relevancia para la valoración de la calidad del proceso, que se
aborda en la sección \ref{sec:critico}.

% ---------------------------------------------------------------------
\subsection{Alcance de la auditoría}

El alcance temporal de la evaluación se extendió desde el momento en que Glamis
Gold asumió la condición de operadora hasta la fecha de emisión del informe, en
mayo de 2010, e incorporó además una mirada prospectiva sobre el cierre y el
post-cierre de la operación. El alcance espacial quedó circunscrito al área de la
licencia de explotación \parencite{oncommonground2010hria}.

Junto al alcance, el informe declaró expresamente una \textbf{limitación}: la
participación del municipio de Sipacapa y de los sectores opositores a la mina
resultó insuficiente, por lo que los propios auditores calificaron sus hallazgos
sobre impactos como parciales.

El registro de esta limitación exige una precisión conceptual que enlaza
directamente con la sección \ref{sec:marco}. En la lógica de la ISO 19011:2018, el
alcance y sus exclusiones no son un preámbulo administrativo: condicionan la
validez de las conclusiones. Si la auditoría se sustenta en un enfoque basado en
la evidencia, y la evidencia se obtiene mediante muestreo, entonces la
representatividad de la muestra determina hasta dónde alcanzan las conclusiones
\parencite{iso2018iso19011}. Una evaluación de impactos sociales que no logra
incorporar el testimonio de la población más crítica no obtiene una muestra
incompleta en sentido meramente cuantitativo: obtiene una muestra sesgada en la
dirección precisa que importa. Cuál sea el peso de ese sesgo se discute en la
sección \ref{sec:critico}; aquí basta con dejar constancia de que los auditores lo
declararon ellos mismos, lo que constituye, en sí, una práctica correcta.

% ---------------------------------------------------------------------
\subsection{Criterios de auditoría aplicados}

Los criterios de auditoría son la vara de medir: el conjunto de requisitos frente
a los cuales se compara la conducta observada. Su elección determina qué puede
llegar a constituir un hallazgo y qué queda, por definición, fuera del examen. La
Tabla \ref{tab:criterios} sistematiza los criterios aplicados en esta evaluación,
distinguiendo su naturaleza jurídica.

% --- TABLA: Criterios de auditoría -----------------------------------
\begin{table}[H]
  \centering
  \caption{Criterios de auditoría aplicados en la evaluación de la Mina Marlin}
  \label{tab:criterios}
  \begin{small}
  \begin{tabularx}{\textwidth}{@{}L{5.2cm} L{3.2cm} Y@{}}
    \toprule
    \textbf{Criterio} & \textbf{Naturaleza} & \textbf{Ámbito que gobierna} \\
    \midrule
    Declaración Universal de Derechos Humanos & Marco internacional &
      Referencia general de derechos. \\
    \addlinespace
    Pacto Internacional de Derechos Civiles y Políticos & Tratado vinculante &
      Seguridad personal, libertad de asociación. \\
    \addlinespace
    Pacto Internacional de Derechos Económicos, Sociales y Culturales &
      Tratado vinculante & Agua, salud, trabajo, medios de vida. \\
    \addlinespace
    Convenio 169 de la {OIT} & Tratado vinculante para el Estado &
      Consulta previa, libre e informada; tierras. \\
    \addlinespace
    Marco «Proteger, Respetar, Remediar» (Ruggie) & Marco de referencia &
      Estructura general del análisis. \\
    \addlinespace
    Herramienta {HRCA} del {DIHR} & Herramienta metodológica &
      Instrumento de verificación, adaptado al caso. \\
    \addlinespace
    {IFC} Performance Standards & Obligación contractual &
      Exigibles por el préstamo de la {IFC}. \\
    \addlinespace
    {ICMM} Mining Principles & Estándar voluntario sectorial &
      Desempeño corporativo. \\
    \addlinespace
    Principios Voluntarios en Seguridad y Derechos Humanos & Estándar voluntario &
      Seguridad privada y pública. \\
    \bottomrule
  \end{tabularx}
  \end{small}
  \par\vspace{3pt}
  \raggedright\footnotesize
  Nota. Elaboración propia a partir de \textcite{oncommonground2010hria}.
\end{table}

La columna central de la tabla revela un rasgo estructural del caso. Los criterios
aplicados pertenecen a tres categorías de muy distinta exigibilidad. Los tratados
internacionales de derechos humanos y el Convenio 169 vinculan al \textbf{Estado
guatemalteco}, no directamente a la empresa. Los IFC Performance Standards
vinculan a la empresa, pero por vía \textbf{contractual}, como condición del
préstamo \parencite{ifcsfmarlin}. Los estándares del ICMM y los Principios
Voluntarios son de adhesión \textbf{voluntaria}. Ninguno de ellos es, en sentido
estricto, una norma sancionable directamente frente al operador por un órgano
administrativo, como sí lo sería el marco peruano expuesto en la sección
\ref{sec:marco}. Esta configuración explica por qué el resultado del proceso
fueron recomendaciones y no medidas correctivas exigibles.

Corresponde señalar, además, una diferencia de fondo respecto de la práctica
auditora convencional. Los criterios empleados son criterios de
\textbf{desempeño} en derechos humanos, no requisitos de un \textbf{sistema de
gestión}. La evaluación no fue una auditoría de certificación conforme a la
ISO 14001:2015 ni pretendió serlo. La consecuencia práctica es que sus resultados
no se formularon como no conformidades en el sentido técnico de la norma, lo que
obliga a este informe a construir una clasificación propia para el análisis de la
sección \ref{sec:hallazgos}. Explicitar esa diferencia evita atribuir a la
evaluación una naturaleza que no tuvo.

% ---------------------------------------------------------------------
\subsection{Metodología y cronograma}

La evaluación se desarrolló entre octubre de 2008 y mayo de 2010, esto es, a lo
largo de aproximadamente diecinueve meses, y se articuló en cinco líneas de
trabajo \parencite{oncommonground2010hria}:

\begin{enumerate}
  \item \textbf{Revisión de políticas, procedimientos y prácticas} de la empresa,
        contrastando el marco formal declarado con su aplicación efectiva.
  \item \textbf{Análisis de datos secundarios y consulta a fuentes expertas},
        incorporando información previamente producida por terceros.
  \item \textbf{Entrevistas individuales y grupales} con múltiples grupos de
        interés: trabajadores, autoridades, comunidades, organizaciones sociales y
        personal de la empresa.
  \item \textbf{Adaptación de la herramienta Human Rights Compliance Assessment}
        del Danish Institute for Human Rights, ajustada a las circunstancias del
        caso.
  \item \textbf{Revisión técnica ambiental independiente}, encargada por el propio
        equipo dentro del proceso de evaluación.
\end{enumerate}

La quinta línea merece un comentario específico, porque es la que convierte a este
proceso en el caso idóneo para el objeto de este informe. Un equipo de perfil
social y jurídico, consciente de la brecha de competencia técnica señalada
anteriormente, subcontrató la pericia ambiental en lugar de eludir la materia o de
pronunciarse sobre ella sin respaldo. Esa decisión metodológica es la que hace de
la evaluación un ejercicio efectivamente \textbf{socioambiental} y no únicamente
social, y es coherente con lo que la ISO 19011:2018 prevé cuando el equipo
auditor requiere conocimientos especializados de los que carece
\parencite{iso2018iso19011}.

\subsubsection{Contraste con las fases de la ISO 19011:2018}

Resulta ilustrativo mapear el proceso efectivamente seguido contra la secuencia de
fases que establece la norma de referencia, expuesta en la sección
\ref{sec:marco}. El ejercicio permite identificar dónde el proceso fue sólido y
dónde quedó débil.

El \textbf{inicio de la auditoría} está documentado con inusual precisión: existe
un memorando de entendimiento escrito, un órgano de gobierno constituido, un
procedimiento de selección del auditor y un conjunto de principios pactados. Es,
probablemente, la fase mejor resuelta del proceso.

La \textbf{revisión de la documentación} y la \textbf{preparación} se corresponden
con las líneas de trabajo primera y segunda, y la adaptación de la herramienta
HRCA equivale a la elaboración de los documentos de trabajo que la norma prevé.

Las \textbf{actividades de recolección de evidencia} se ejecutaron mediante
entrevistas, análisis documental y la revisión técnica encargada, cubriendo las
tres modalidades que la norma contempla. Es en esta fase, sin embargo, donde se
localiza la debilidad declarada: el muestreo de personas entrevistadas no alcanzó
la cobertura prevista respecto de Sipacapa y de los sectores opositores.

El \textbf{informe} se elaboró y se publicó íntegramente, en versión bilingüe, lo
que satisface con holgura el requisito de comunicación de resultados.

El \textbf{seguimiento} constituye el punto más problemático del contraste. La
norma concibe la auditoría como un ciclo que se cierra verificando la
implementación efectiva de las acciones correctivas. Aquí la empresa se comprometió
a emitir una respuesta pública y un plan de acción, con actualizaciones
publicadas en 2010, 2011 y 2017 \parencite{oncommonground2010hria,
goldcorp2023marlin}, pero la verificación de ese cumplimiento quedó en manos de la
propia empresa, sin que el equipo auditor original conservara mandato para
comprobarlo. La cuestión de si las no conformidades llegaron efectivamente a
cerrarse se retoma en las secciones \ref{sec:hallazgos} y \ref{sec:critico}.

% --- FIGURA 3: Línea de tiempo (dibujada en TikZ) --------------------
\begin{figure}[H]
  \centering
  \begin{tikzpicture}[
      hito/.style={circle,fill=azuloscuro,inner sep=1.4pt},
      arr/.style={-{Stealth[length=2.5mm]},very thick,azuloscuro},
      anio/.style={font=\scriptsize\bfseries,anchor=north,inner sep=1.5pt},
      arriba/.style={font=\scriptsize,align=center,anchor=south,inner sep=2pt},
      abajo/.style={font=\scriptsize,align=center,anchor=north,inner sep=2pt}
    ]
    % Eje
    \draw[arr] (-0.4,0) -- (14.2,0);

    % Hitos y años, pegados al eje
    \foreach \x/\y in {0/2004, 1.3/2005, 2.6/2006, 4.3/2008, 5.5/2009,
                       6.7/2010, 8.3/2011, 13.4/2017} {
      \node[hito] at (\x,0) {};
      \node[anio] at (\x,-0.14) {\y};
    }

    % Etiquetas por encima del eje
    \node[arriba] at (0,0.14)   {Revisión\\independiente\\del EIA};
    \node[arriba] at (2.6,0.14) {Cierre del\\caso CAO};
    \node[arriba] at (5.5,0.14) {Trabajo\\de campo};
    \node[arriba] at (8.3,0.14) {La CIDH retira\\la solicitud\\de suspensión};

    % Etiquetas por debajo, dejando sitio a la fila de años
    \node[abajo] at (1.3,-0.62)  {Consulta de\\Sipacapa;\\queja ante\\la CAO};
    \node[abajo] at (4.3,-0.62)  {MOU (mar.)\\Selección de\\la auditora\\(oct.)};
    \node[abajo] at (6.7,-0.62)  {\textbf{Publicación}\\\textbf{de la HRIA}\\(mayo)\\Medidas CIDH\\OIT/CEACR};
    \node[abajo] at (13.4,-0.62) {Cierre de\\la mina\\(31 mayo)};

    % Duración del proceso de auditoría
    \draw[azuloscuro,thick,|-|] (4.3,1.55) -- (6.7,1.55)
      node[midway,above,font=\scriptsize] {Proceso de auditoría (19 meses)};
  \end{tikzpicture}
  \caption{Cronología del proceso de auditoría de la Mina Marlin (2004--2017)}
  \label{fig:metodologia}
  \par\vspace{2pt}
  \raggedright\footnotesize
  Nota. El eje no es estrictamente proporcional entre 2011 y 2017.
  Elaboración propia a partir de \textcite{oncommonground2010hria},
  \textcite{cao2006marlin} y \textcite{cidh2010medidas}.
\end{figure}
